\documentclass[11pt,a4paper]{altacv}

\geometry{left=1.5cm,right=1.5cm,top=1.cm,bottom=1.2cm,columnsep=1.5cm}

\usepackage[utf8]{inputenc}
\usepackage[T1]{fontenc}
\usepackage[french]{babel}
\usepackage{paracol}
\usepackage{xcolor}
\usepackage{hyperref}
\usepackage{fontawesome5}
\usepackage{setspace}

\definecolor{SlateGrey}{HTML}{2E2E2E}
\definecolor{LightGrey}{HTML}{666666}
\definecolor{DarkPastelRed}{HTML}{8AC0D9}
\definecolor{PastelRed}{HTML}{00B0DA}
\definecolor{GoldenEarth}{HTML}{B4AD0C}
\colorlet{name}{black}
\colorlet{tagline}{PastelRed}
\colorlet{heading}{DarkPastelRed}
\colorlet{headingrule}{GoldenEarth}
\colorlet{subheading}{PastelRed}
\colorlet{accent}{PastelRed}
\colorlet{emphasis}{SlateGrey}
\colorlet{body}{LightGrey}

\renewcommand{\familydefault}{\sfdefault}

\name{\Large Guillaume BOILEAU}
\tagline{Ingénieur AIT / IVVQ spatial | Intégration, validation fonctionnelle, analyse d'essais et anomalies techniques}
\personalinfo{
  \email{guillaume.boileaupro@gmail.com}
  \phone{+33 6 59 94 85 84}
  \location{Alpes-Maritimes, France}
  \linkedin{guillaume-boileau-phd-891b81265}
}

\setlength{\parskip}{0.75\baselineskip}

\begin{document}
\makecvheader
\vspace{-0.6cm}

\cvsection{Lettre de motivation}
\vspace{-0.4cm}

{\color{accent}\textbf{Objet :}} \textit{Candidature -- Ingénieur AIT Spatial -- Cannes}

{\color{body}

Madame, Monsieur,

Je vous adresse ma candidature pour le poste d'\textbf{Ingénieur AIT Spatial} à Cannes. Votre offre m'intéresse particulièrement car elle combine plusieurs dimensions directement cohérentes avec mon parcours : environnement spatial, intégration de systèmes complexes, validation fonctionnelle, analyse de résultats d'essais, traitement d'anomalies techniques, rédaction de procédures et interface avec des équipes pluridisciplinaires.

Mon parcours s'est construit à l'interface entre systèmes spatiaux, validation, traitement du signal, analyse de données, développement logiciel et coordination technique. J'ai travaillé dans des environnements exigeants où la robustesse des méthodes, la qualité des résultats, la traçabilité des hypothèses, la documentation technique et l'analyse d'écarts entre comportement attendu et comportement observé sont essentielles.

Au \textbf{CNES}, j'ai contribué à des activités de développement, d'intégration, de comparaison et de validation pour la mission spatiale \textbf{LISA}, mission ESA/NASA de grande envergure. J'y ai notamment développé \textbf{CATpyLISA}, une plateforme Python destinée à l'intégration, la comparaison et la validation de modèles et de chaînes d'analyse complexes. Ce travail m'a amené à structurer des workflows techniques, à organiser des campagnes de comparaison, à analyser des incohérences, à documenter les résultats et à produire des éléments de synthèse exploitables par différents contributeurs.

Cette expérience est directement transposable à une logique \textbf{AIT / IVVQ} : préparer des activités de validation, définir des critères d'analyse, vérifier la cohérence des résultats, investiguer les écarts, produire une preuve documentée et assurer l'interface avec les équipes impliquées. Même si mon parcours n'est pas celui d'un ingénieur AIT satellite classique issu directement des bancs RF ou EGSE/MGSE, il m'a donné une solide pratique des environnements spatiaux, de la validation système, de l'analyse de performance, du traitement d'anomalies et de la documentation technique.

Je dispose également d'un socle pertinent pour aborder les essais fonctionnels et RF : traitement du signal, analyse spectrale, bruit instrumental, environnements bas signal-sur-bruit, propagation de signaux, analyse de mesures et identification de limitations de performance. Mes travaux au sein des collaborations LISA, LIGO--Virgo--KAGRA et Einstein Telescope m'ont habitué à raisonner sur des systèmes complexes, des instruments de haute précision, des sources de bruit, des contraintes de mesure et des effets pouvant limiter la performance globale d'un système.

Plus récemment, chez \textbf{VEV Platform Services France}, j'ai travaillé sur des services logiciels liés aux infrastructures de recharge électrique, avec des interfaces entre plateforme logicielle, équipements connectés, flux de données opérationnels et contraintes terrain. Cette expérience industrielle m'a permis de renforcer ma capacité à analyser des incidents en environnement opérationnel, à suivre des interfaces, à diagnostiquer des comportements anormaux, à documenter des observations techniques et à contribuer à la fiabilité de systèmes connectés.

Pour ALTEN, je peux apporter un profil hybride \textbf{spatial, système, validation, logiciel et analyse technique}. Je suis capable de comprendre rapidement un environnement complexe, de structurer une démarche d'essai ou de validation, d'analyser des résultats, d'identifier des écarts, de documenter des anomalies et de dialoguer avec des équipes conception, qualité, méthodes et programme. Je suis également habitué à travailler en anglais technique dans des collaborations internationales et je suis basé dans les Alpes-Maritimes, ce qui me permet de m'inscrire naturellement dans l'écosystème spatial de Cannes.

Je souhaite aujourd'hui mettre cette double culture, issue du spatial scientifique et du logiciel industriel, au service de programmes satellites concrets. Le poste proposé par ALTEN représente pour moi une continuité logique vers des activités d'intégration, d'essais, de validation et de support technique au plus près des systèmes spatiaux.

Je serais heureux de pouvoir échanger avec vous afin de vous présenter plus en détail mon parcours et la manière dont mon expérience en validation de systèmes complexes, analyse de performances, traitement d'anomalies et coordination technique peut contribuer efficacement à vos activités AIT spatiales.

Je vous prie d'agréer, Madame, Monsieur, l'expression de mes salutations distinguées.

\textbf{Guillaume Boileau}

}

\end{document}
